On both datasets, the proposed pipeline achieved the highest session-mean $F_1$, followed by \rev{the BLINKER-style baseline} and then MNE. On Internal (Table~\ref{tab:f1_significance}), the proposed pipeline reached 88.25\%, compared with 76.54\% for \rev{the BLINKER-style baseline} and 66.64\% for MNE. Its advantage was 11.7 percentage points over \rev{the BLINKER-style baseline} and 21.6 percentage points over MNE, and both differences were significant (adjusted $p < 0.05$). On Cao2018 (Table~\ref{tab:f1_significance}), the proposed pipeline likewise reached 80.68\%, against 69.40\% for \rev{the BLINKER-style baseline} and 56.27\% for MNE. The margins were 11.3 and 24.4 percentage points, and both were again significant (adjusted $p < 0.05$). All three methods scored lower on Cao2018 than on Internal, but the decline was not uniform. The proposed pipeline and \rev{the BLINKER-style baseline} lost 7.6 and 7.1 percentage points, respectively, so the margin between them stayed close to 11 percentage points on both datasets. MNE, in contrast, lost 10.4 percentage points, which widened its gap to the proposed pipeline from 21.6 to 24.4 percentage points and its gap to \rev{the BLINKER-style baseline} from 9.9 to 13.1 percentage points.
